\subsection{Distinguished subgroupoids; quotients of groupoids}

DEFINITION 7. --- Let G be a groupoid. One says that a subgroupoid H of G is distinguished if Som(H) = Som(G) and if, for every pair (a, b) of vertices of H and every arrow \(f \in \mathrm{G}_{a,b}\) , one has \(Int(f)(H_b) = H_a\) .

Let G be a groupoid and let H be a subgroupoid of G whose set of vertices is equal to Som(G). To verify that H is distinguished, it suffices to show that \(fH_{b}f^{-1} \subset H_{a}\) for all \(a \in \text{Som}(G)\) , all \(b \in \text{Som}(\mathbf{G})\) and every \(f \in \mathbf{G}_{a,b}\) . Indeed, if this is so, we also have \(f^{-1}\mathrm{H}_a f \subset \mathrm{H}_b\) for every arrow \(f \in \mathbf{G}_{b,a}\) , that is, \(\mathrm{H}_a \subset f\mathrm{H}_b f^{-1}\) , and, consequently, \(f\mathrm{H}_b f^{-1} = \mathrm{H}_a\) .

Let G be a groupoid and let H be a distinguished subgroupoid of G. For every vertex a of G, the subgroup \(H_{a}\) of \(G_{a}\) is a distinguished subgroup of \(G_{a}\) .

Let \(\varphi\colon\mathrm{G}\to\mathrm{G}^{\prime}\) be a morphism of groupoids. The kernel of \(\varphi\) (II, p. 166, example 5) is a distinguished subgroupoid of G. Indeed, let \(f\) be an arrow in G connecting \(a\) to \(b\) and let \(g\in\mathrm{Ker}(\varphi)_{b}\) ; we then have \(\varphi(fgf^{-1})=\varphi(f)\varphi(g)\varphi(f^{-1})=\varphi(f)e_{\varphi(b)}\varphi(f)^{-1}=e_{\varphi(a)}\) , whence the inclusion \(f\mathrm{Ker}(\varphi)_{b}f^{-1}\subset\mathrm{Ker}(\varphi)_{a}\) .

Let G be a groupoid. The groupoid G and the subgroupoid of G whose set of arrows is the set of identity elements of G are distinguished subgroupoids of G. The intersection of a family of distinguished subgroupoids of G is a distinguished subgroupoid of G. In particular, for every subset F ⊂ Fl(G), there exists a smallest distinguished subgroupoid of G whose set of arrows contains F. It is called the distinguished subgroupoid generated by F.

Let H be a distinguished subgroupoid of G. Let R be the equivalence relation in \(\mathrm{Fl}(\mathrm{G})\) defined by \(R\{f,g\}\) if and only if there exist arrows x and y in \(\mathrm{Fl}(\mathrm{H})\) such that f = xgy. If f and g are arrows of G equivalent modulo R, their origins (resp. their targets) belong to the same orbit of H. Set \(F' = \mathrm{Fl}(\mathrm{G}) / \mathcal{R}\) and denote by \(o'\) and \(t'\) the maps from \(F'\) to \(\mathrm{Orb}(\mathrm{H})\) deduced from o and t by passing to quotients. Denote by G/H the quiver \((\mathrm{Orb}(\mathrm{H}), \mathrm{F}', o', t')\) .

Lemma. --- Given composable arrows u and v of G/H, one can choose representatives f and g of u and v in Fl(G) which are composable. The class modulo R of the product fg does not depend on the choice of f and g.

Let \(f\) and \(g\) be arbitrary representatives of \(u\) and \(v\) in \(\mathrm{Fl(G)}\) . The target of \(f\) and the origin of \(g\) belong to the same orbit of H, hence can be connected by an arrow \(x\) of H (II, p. 162, prop. 1). The arrows \(fx\) and \(g\) are then representatives of \(u\) and \(v\) in F that are composable in G. This proves the first assertion.

Now let \(f\) , \(g\) on the one hand, and \(f'\) , \(g'\) on the other hand, be representatives of \(u\) and \(v\) which are composable in G. By hypothesis, there exist arrows \(x\) , \(y\) , \(z\) , \(t\) in H such that one has \(f' = xfy\) and \(g' = zgt\) . One then has \(yz \in \mathrm{H}_{t(f)}\) and \(f'g' = x f y z g t = x f(yz)f^{-1}f g t\) . Since H is a distinguished subgroupoid of G, the arrow \(f(yz)f^{-1}\) is an arrow of H. The same is therefore true of the arrow \(x f(yz)f^{-1}\) , which proves that \(f'g'\) and \(fg\) are equivalent modulo \(\mathcal{R}\) .

By virtue of this lemma, there exists a unique composition law \(m\) in the quiver G/H such that, for every pair \((u,v)\) of composable arrows, \(m(u,v)\) is the class modulo \(\mathcal{R}\) of the product \(fg\) , for every pair \((f,g)\) of composable arrows of G such that \(u\) is the class of \(f\) and \(v\) that of \(g\) . Endowed with this composition law, G/H is a groupoid. Let \(p_1\colon \operatorname{Som}(\mathrm{G}) \to \operatorname{Som}(\mathrm{G}/\mathrm{H})\) and \(p_2\colon \operatorname{Fl}(\mathrm{G}) \to \operatorname{Fl}(\mathrm{G}/\mathrm{H})\) be the canonical surjections. The pair \(p = (p_1,p_2)\) is a groupoid morphism from G to G/H whose kernel is the distinguished subgroupoid H.

DEFINITION 8. --- We say that G/H is the quotient groupoid of G by H and that p: G → G/H is the canonical morphism.

Remarks. --- 1) The map Som(p) defines, by passing to quotients, a bijection from the set of orbits of G to the set of orbits of G/H. In particular, G is transitive if and only if G/H is.

\begin{enumerate}
\def\labelenumi{\arabic{enumi})}
\setcounter{enumi}{1}
\tightlist
\item
  Let \(a\) and \(b\) be vertices of G. The map from \(\mathrm{G}_{a,b}\) to \((\mathrm{G} / \mathrm{H})_{p(a),p(b)}\) deduced from \(p\) is surjective. Indeed, let \(u\) be an arrow from G/H connecting \(p(a)\) to \(p(b)\) ; it is the class modulo \(\mathcal{R}\) of an arrow \(f\) of G. The origin \(o(f)\) of \(f\) belongs to the orbit of \(a\) in H; there therefore exists an arrow \(x\) in H connecting \(a\) to \(o(f)\) . Similarly, there exists an arrow \(y\) in H that connects \(t(f)\) to \(b\) . Then, \(f' = xfy\) is an element of \(\mathrm{G}_{a,b}\) whose class modulo \(\mathcal{R}\) is \(u\) .
\end{enumerate}

PROPOSITION 2. --- Let a be a vertex of G. The group homomorphism \(p_{a}\colon\mathrm{G}_{a}\to(\mathrm{G}/\mathrm{H})_{p(a)}\) deduced from p by passing to isotropy groups is surjective; its kernel is \(\mathrm{H}_{a}\) .

The homomorphism \(p_a\) is surjective by virtue of the preceding remark. Its kernel contains \(\mathrm{H}_a\) by construction of the groupoid \(\mathrm{G / H}\) . Let \(f\) be an element of \(\mathrm{G}_a\) such that \(p_a(f) = e_{p(a)}\) . This means that \(f\) and \(e_a\) are equivalent modulo \(\mathcal{R}\) ; there then exist arrows \(x\) and \(y\) of H such that \(f = xe_a y\) . Necessarily, \(x\) and \(y\) belong to \(\mathrm{H}_a\) , whence \(f \in \mathrm{H}_a\) .

PROPOSITION 3. --- Let G be a groupoid, let H be a distinguished subgroupoid of G, and let p: G → G/H be the canonical morphism. Let \(\varphi\colon G\to G'\) be a morphism of groupoids such that \(H\subset\operatorname{Ker}(\varphi)\) . There exists a unique groupoid morphism \(\overline{\varphi}\colon G/H\to G'\) such that \(\overline{\varphi}\circ p=\varphi\) .

We say that \(\overline{\varphi}\) is the groupoid morphism induced by \(\varphi\) by passing to the quotient.

The uniqueness of such a morphism is evident, since the maps \(\operatorname{Som}(p)\) and \(\operatorname{Fl}(p)\) are surjective.

Let \(a\) and \(b\) be vertices of G. If they are in the same orbit of H, there exists an arrow \(f\) connecting \(a\) to \(b\) in H and one has \(\varphi(f) = e_{\varphi(a)}\) . In particular, \(\varphi(a) = \varphi(b)\) . Consequently, the map \(\operatorname{Som}(\varphi)\) defines, by passing to the quotient, a map \(\overline{\varphi}_1\colon \operatorname{Orb}(\mathrm{H}) \to \operatorname{Som}(\mathrm{G}')\) . Let \(f\) and \(g\) be arrows in G. If \(f\) and \(g\) are equivalent modulo \(\mathcal{R}\) , there exist arrows \(x\) and \(y\) in H such that \(f = xgy\) . Consequently, \(\varphi(f) = \varphi(x)\varphi(g)\varphi(y) = \varphi(g)\) since \(\varphi(x)\) and \(\varphi(y)\) are identity elements. The map \(\operatorname{Fl}(\varphi)\) therefore defines, by passing to the quotient, a map \(\overline{\varphi}_2\colon \operatorname{Fl}(\mathrm{G}) / \mathcal{R} \to \operatorname{Fl}(\mathrm{G}')\) .

Let \(f\) and \(g\) be arrows of \(G\) which are composable; denote by \(u\) and \(v\) their classes in \(\mathrm{Fl}(\mathrm{G} / \mathrm{H})\) . We have \(\overline{\varphi}_2(uv) = \varphi(fg) = \varphi(f)\varphi(g) = \overline{\varphi}_2(u)\overline{\varphi}_2(v)\) .

The pair \(\overline{\varphi} = (\overline{\varphi}_1, \overline{\varphi}_2)\) is a morphism of groupoids from G/H to G', and one has \(\overline{\varphi} \circ p = \varphi\) .

